% GENERATED FILE. Edit canonical lessons, then run npm run generate:books.
% canonical-source-hash: fnv1a64:261f9a5cf13c7c36
% canonical-lessons: ES-C574-enterarse, ES-C574-recibir, ES-C574-notificacion, ES-C574-telediario, ES-C574-incidente

\chapter{Cómo me enteré}
\label{ch:es-como-me-entere}
\hlchaptermodality{574}

\begin{chapteropening}
\textbf{By the end of this chapter:} \emph{I can say how I found out about an incident: I got a notification or saw it on the TV news.}

Complete the last lesson of chapter 574: I can say how I found out about an incident: I got a notification or saw it on the TV news.
\end{chapteropening}

\section[enterarse]{enterarse — to hear about, to find out}

\label{lesson:ES-C574-enterarse}



\begin{quote}
Before the new one: say the Spanish for to suffer, then the Spanish for to turn out.
\end{quote}

\subsection*{You'll want to know: enterarse}
\textbf{enterarse} — \textquotedblleft{}to hear about, to find out\textquotedblright{}.

\begin{quote}
Me enteré de la noticia por la radio.
\end{quote}

\begin{quote}
\emph{I heard the news on the radio.}
\end{quote}

\emph{Enterarse de algo} — to find out about something. \emph{¿Te has enterado?} — have you heard?

\begin{cousinweb}
\emph{enterarse} is built on \emph{entero} 'whole', inherited from Latin \emph{integer} 'untouched, whole' (Lewis \& Short). To find out is to get the whole story.
\end{cousinweb}

\begin{grammarlens}[title={What you've built}]
One more word for telling, explaining and describing at B1.
\end{grammarlens}

\subsection*{Your turn}
\begin{itemize}
\raggedright
  \item \emph{Say it:} \emph{enterarse}
  \item \emph{Say it:} \emph{enterarse}, once more
  \item \emph{Say it:} a sentence of your own with \emph{enterarse}
  \item \emph{Recall:} say the Spanish for to suffer, then the Spanish for to turn out, then say \emph{enterarse} again
\end{itemize}

\begin{tcolorbox}[breakable,colback=teal!4,colframe=teal!35!black,title={Before you move on}]
What does enterarse mean? (\textquotedblleft{}To hear about, to find out\textquotedblright{}.) Say it once more.
\end{tcolorbox}
\section[recibir]{recibir — to receive, to get}

\label{lesson:ES-C574-recibir}



\begin{quote}
Before the new one: say the Spanish for to turn out, then the Spanish for to hear about.
\end{quote}

\subsection*{You'll want to know: recibir}
\textbf{recibir} — \textquotedblleft{}to receive, to get\textquotedblright{}.

\begin{quote}
Recibí una carta de mi tía.
\end{quote}

\begin{quote}
\emph{I got a letter from my aunt.}
\end{quote}

It also means to welcome: \emph{Nos recibieron muy bien.}

\begin{cousinweb}
\emph{recibir} is inherited from Latin \emph{recipere} 'to take back, to receive' (Lewis \& Short), built on \emph{re-} and \emph{capere} 'to take'.
\end{cousinweb}

\begin{grammarlens}[title={What you've built}]
One more word for telling, explaining and describing at B1.
\end{grammarlens}

\subsection*{Your turn}
\begin{itemize}
\raggedright
  \item \emph{Say it:} \emph{recibir}
  \item \emph{Say it:} \emph{recibir}, once more
  \item \emph{Say it:} a sentence of your own with \emph{recibir}
  \item \emph{Recall:} say the Spanish for to turn out, then the Spanish for to hear about, then say \emph{recibir} again
\end{itemize}

\begin{tcolorbox}[breakable,colback=teal!4,colframe=teal!35!black,title={Before you move on}]
What does recibir mean? (\textquotedblleft{}To receive, to get\textquotedblright{}.) Say it once more.
\end{tcolorbox}
\section[la notificación]{la notificación — the notification, the alert}

\label{lesson:ES-C574-notificacion}



\begin{quote}
Before the new one: say the Spanish for to hear about, then the Spanish for to receive.
\end{quote}

\subsection*{You'll want to know: la notificación}
\textbf{la notificación} — \textquotedblleft{}the notification, the alert\textquotedblright{}.

\begin{quote}
Me llegó una notificación al móvil.
\end{quote}

\begin{quote}
\emph{I got a notification on my phone.}
\end{quote}

Also an official notice: \emph{una notificación del juzgado}.

\begin{cousinweb}
\emph{notificación} is built on \emph{notificar}, a learned borrowing from Latin \emph{nōtificāre} 'to make known' (Lewis \& Short), built on \emph{nōtus} 'known' and \emph{facere} 'to make'.
\end{cousinweb}

\begin{grammarlens}[title={What you've built}]
One more word for telling, explaining and describing at B1.
\end{grammarlens}

\subsection*{Your turn}
\begin{itemize}
\raggedright
  \item \emph{Say it:} \emph{la notificación}
  \item \emph{Say it:} \emph{la notificación}, once more
  \item \emph{Say it:} a sentence of your own with \emph{la notificación}
  \item \emph{Recall:} say the Spanish for to hear about, then the Spanish for to receive, then say \emph{la notificación} again
\end{itemize}

\begin{tcolorbox}[breakable,colback=teal!4,colframe=teal!35!black,title={Before you move on}]
What does la notificación mean? (\textquotedblleft{}The notification, the alert\textquotedblright{}.) Say it once more.
\end{tcolorbox}
\section[el telediario]{el telediario — the television news}

\label{lesson:ES-C574-telediario}



\begin{quote}
Before the new one: say the Spanish for to receive, then the Spanish for the notification.
\end{quote}

\subsection*{You'll want to know: el telediario}
\textbf{el telediario} — \textquotedblleft{}the television news\textquotedblright{}.

\begin{quote}
Vi la noticia en el telediario de las nueve.
\end{quote}

\begin{quote}
\emph{I saw the story on the nine o'clock news.}
\end{quote}

The word in Spain; in Latin America people say \emph{el noticiero}.

\begin{cousinweb}
\emph{telediario} is built on \emph{tele-}, from Greek \emph{tēle} 'far', and \emph{diario} 'daily', which comes from Latin \emph{diārium} 'a daily allowance' (Lewis \& Short), from \emph{diēs} 'day'.
\end{cousinweb}

\begin{grammarlens}[title={What you've built}]
One more word for telling, explaining and describing at B1.
\end{grammarlens}

\subsection*{Your turn}
\begin{itemize}
\raggedright
  \item \emph{Say it:} \emph{el telediario}
  \item \emph{Say it:} \emph{el telediario}, once more
  \item \emph{Say it:} a sentence of your own with \emph{el telediario}
  \item \emph{Recall:} say the Spanish for to receive, then the Spanish for the notification, then say \emph{el telediario} again
\end{itemize}

\begin{tcolorbox}[breakable,colback=teal!4,colframe=teal!35!black,title={Before you move on}]
What does el telediario mean? (\textquotedblleft{}The television news\textquotedblright{}.) Say it once more.
\end{tcolorbox}
\section[el incidente]{el incidente — the incident}

\label{lesson:ES-C574-incidente}



\begin{quote}
Before the new one: say the Spanish for the notification, then the Spanish for the television news.
\end{quote}

\subsection*{You'll want to know: el incidente}
\textbf{el incidente} — \textquotedblleft{}the incident\textquotedblright{}.

\begin{quote}
El vuelo llegó sin incidentes.
\end{quote}

\begin{quote}
\emph{The flight arrived without incident.}
\end{quote}

\begin{cousinweb}
\emph{incidente} is a learned borrowing from Latin \emph{incidēns} 'falling upon', from \emph{incidere} 'to fall into, to happen' (Lewis \& Short), built on \emph{in} and \emph{cadere} 'to fall'.
\end{cousinweb}

\begin{grammarlens}[title={What you've built}]
One more word for telling, explaining and describing at B1.
\end{grammarlens}

\subsection*{Your turn}
\begin{itemize}
\raggedright
  \item \emph{Say it:} \emph{el incidente}
  \item \emph{Say it:} \emph{el incidente}, once more
  \item \emph{Say it:} a sentence of your own with \emph{el incidente}
  \item \emph{Recall:} say the Spanish for the notification, then the Spanish for the television news, then say \emph{el incidente} again
\end{itemize}

\begin{tcolorbox}[breakable,colback=teal!4,colframe=teal!35!black,title={Before you move on}]
What does el incidente mean? (\textquotedblleft{}The incident\textquotedblright{}.) Say it once more.
\end{tcolorbox}
